\documentclass[reqno]{amsart}

\theoremstyle{plain}
\newtheorem{theorem}{Theorem}
\newtheorem{lemma}[theorem]{Lemma}
\newtheorem{proposition}[theorem]{Proposition}
\theoremstyle{definition}

\newtheorem{definition}[theorem]{Definition}

\begin{document}

\author{Paul D. Nelson}

\section{Question}
Let \(F\) be a non-archimedean local field with ring of integers \(\mathfrak o\).  Let $N_r$ denote the subgroup of $\mathrm{GL}_{r}(F)$ consisting of upper-triangular unipotent elements.  Let \(\psi:F\to \mathbb C^\times\) be a nontrivial additive character of conductor \(\mathfrak o\), identified in the standard way with a generic character of $N_r$.  Let \(\Pi\) be a generic irreducible admissible representation of \(\mathrm{GL}_{n + 1}(F)\), realized in its \(\psi^{-1}\)-Whittaker model \(\mathcal W(\Pi,\psi^{-1})\).  Must there exist \(W\in \mathcal W(\Pi,\psi^{-1})\) with the following property?

Let $\pi$ be a generic irreducible admissible representation of \(\mathrm{GL}_{n}(F)\), realized in its $\psi$-Whittaker model \(\mathcal W(\pi,\psi)\).  Let $\mathfrak{q}$ denote the conductor ideal of $\pi$, let \(Q\in F^\times\) be a generator of \(\mathfrak q^{-1}\), and set
\[
  u_Q := I_{n+1} + Q\,E_{n,n+1} \in \mathrm{GL}_{n + 1}(F),
\]
where \(E_{i, j}\) is the matrix with a \(1\) in the \((i, j)\)-entry and \(0\) elsewhere.  For some \(V\in \mathcal W(\pi,\psi)\), the local Rankin--Selberg integral
\[
  \int_{N_n\backslash \mathrm{GL}_{n}(F)} W(\operatorname{diag}(g,1) u_Q)\,V(g)\,|\det g|^{s-\frac12}\,dg
\]
is finite and nonzero for all \(s\in\mathbb C\).

\section{Statement}

Let \(F\) be a non-archimedean local field with ring of integers \(\mathfrak o\). Let \(\psi:F\to \mathbb C^\times\) be a nontrivial additive character of conductor \(\mathfrak o\).  We write
\begin{equation*}
  G_r:=\operatorname{GL}_r(F),
\end{equation*}
and let \(N_r<G_r\) denote the subgroup of upper-triangular unipotent elements.  We embed \(G_n \hookrightarrow G_{n + 1}\) as the upper-left block.  We write \(E_{i j}\) for the matrix with a \(1\) in the \((i, j)\)-entry and \(0\) elsewhere.

A more precise form of the following ``lemma'' will appear in forthcoming joint work with Subhajit Jana.  It says informally that pure unipotent translates of fixed vectors in the Whittaker model of a representation of \(G_{n+1}\) may serve as test vectors for Rankin--Selberg integrals against all representations of \(G_n\) with a given conductor.

\begin{theorem}\label{thm:main}
  Let \(\Pi\) be a generic irreducible admissible representation of \(G_{n+1}\), realized in its \(\psi^{-1}\)-Whittaker model \(\mathcal W(\Pi,\psi^{-1})\). Then there exists \(W\in \mathcal W(\Pi,\psi^{-1})\) with the following property.  Let $\pi$ be a generic irreducible admissible representation of \(G_n\), realized in its $\psi$-Whittaker model \(\mathcal W(\pi,\psi)\).  Let $\mathfrak{q}$ denote the conductor ideal of $\pi$, let \(Q\in F^\times\) be a generator of \(\mathfrak q^{-1}\), and set
  \[
    u_Q := I_{n+1} + Q\,E_{n,n+1} \in G_{n + 1}.
  \]
  There exists \(V\in \mathcal W(\pi,\psi)\) so that the local Rankin--Selberg integral
  \[
    \int_{N_n\backslash G_n} W(\operatorname{diag}(g,1) u_Q)\,V(g)\,|\det g|^{s-\frac12}\,dg
  \]
is finite and nonzero for all \(s\in\mathbb C\).
\end{theorem}

\section{Context}

Rankin--Selberg local zeta integrals arise as proportionality factors relating global Rankin--Selberg integrals and $L$-functions.  The above result provides test vectors, obtained via pure translates of fixed vectors, that work simultaneously for all representations of the smaller group having some given conductor.  Such results are sometimes useful in global applications because they relate problems concerning $L$-functions (subconvexity, moment asymptotics, ...) to problems concerning automorphic forms (quantitative equidistribution, ...).  The $n = 1$ case follows from standard properties of Gauss sums and stationary phase analysis in one variable; it has been applied in, e.g., \cite{MR780071, michel-2009}.  For general \(n\), \cite{MR4053059} contains a similar result, but with an average over many unipotent translates rather than just one.

\section{Proof}\label{sec:proof}
We first sketch the argument.  The basic idea is to apply the Godement--Jacquet functional to the Whittaker function on the smaller group.  This is readily seen to relate the unipotent-shifted Rankin--Selberg integral to an integral involving a translate of the standard congruence subgroup $K_1(\mathfrak{q}) \leq \mathrm{GL}_n(\mathfrak{o})$, consisting of matrices whose last row is congruent to \((0,\dotsc,0,1)\) modulo \(\mathfrak q\).  We then conclude via newvector theory.

Turning to details, we recall that $F$ is a non-archimedean local field, with ring of integers $\mathfrak{o}$.  We denote by $\mathfrak{p}$ the maximal ideal and $q$ the residue field cardinality.  We set $K_r := \mathrm{GL}_r(\mathfrak{o})$ and equip $G_r$ and $N_r$ with the Haar measures assigning volume one to $K_r$ and $N_r \cap K_r$, respectively.  As in the theorem statement, we write $\Pi$ (resp.\ $\pi$) for a generic irreducible representation of $G_{n + 1}$ (resp.\ $G_n$).

We continue to denote by $\mathfrak{q}$ the conductor ideal of $\pi$, defined to be the smallest ideal for which $\pi$ has a nonzero vector fixed by $K_1(\mathfrak{q})$.  We choose a generator $Q$ for $\mathfrak{q}^{-1}$, so that $\lvert Q \rvert =[\mathfrak{o} : \mathfrak{q}]$.  We recall (see \cite{MR620708, MR3138844}) that $\lvert Q \rvert$ (and hence $\mathfrak{q}$) may also be characterized in terms of the local $\varepsilon$-factor of $\pi$:
\begin{equation}\label{eq:cuipy1wony}
  \varepsilon(\tfrac{1}{2} + s, \pi, \psi) = \lvert Q \rvert^{- s} \varepsilon(\tfrac{1}{2}, \pi, \psi).
\end{equation}

We recall the functional equation of Godement--Jacquet \cite[Theorem 3.3]{MR0342495}.

\begin{lemma}\label{lemma:cq4itap0v1}
  Let $f$ be a matrix coefficient of $\pi$, and let $\phi \in \mathcal{S}(M_n(F))$.  For $s \in \mathbb{C}$, the local zeta integral
  \begin{equation}\label{eq:cq73uvikdb}
    Z(\phi, f, s) := \int_{G_n}
    \phi(g) f(g) \lvert \det g \rvert^{\frac{n - 1}{2} + s} \, d g,
  \end{equation}
  converges absolutely for $\Re(s)$ sufficiently large.  It extends to a meromorphic function on the complex plane for which the ratio
  \begin{equation*}
    \frac{Z(\phi, f, s)}{L(s, \pi)}
  \end{equation*}
  is holomorphic.  It satisfies the local functional equation
  \begin{equation}\label{eq:cq4itiiq9o}
    \gamma(s, \pi, \psi) Z(\phi, f, s)
    = Z(\phi^\wedge, f^\vee, 1-s),
  \end{equation}
  where
  \begin{equation*}
    \gamma(s, \pi, \psi) = \varepsilon(s, \pi, \psi) \frac{L(1 - s, \tilde{\pi})}{L(s, \pi)},
  \end{equation*}
  with $\tilde{\pi}$ the contragredient of $\pi$, and where the Fourier transform is defined by
  \begin{equation*}
    f^\vee (g) := f(g^{-1}),
  \end{equation*}
  \begin{equation*}
    \phi^\wedge(x) := \int_{M_n(F)}
    \phi(y) \psi(\operatorname{trace}(x y)) \, d y,
  \end{equation*}
  with $M_n$ the space of $n \times n$ matrices and the Haar measure normalized to be self-dual with respect to $\psi$.  Moreover, both of the zeta integrals in \eqref{eq:cq4itiiq9o} converge absolutely provided that, e.g., $\pi$ is unitary and generic and $\Re(s) = 1/2$.
\end{lemma}
We recall that a matrix coefficient of $\pi$ is a linear combination of functions of the form $f(g) = \ell(g v)$, where $v \in \pi$ and $\ell$ lies in the contragredient of $\pi$ (i.e., the admissible dual).  The conclusions of Lemma \ref{lemma:cq4itap0v1} remain valid for more general coefficients of $\pi$.  For instance, suppose more generally that $f$ is of the same form, but with $\ell$ allowed to be any linear functional on $\pi$ (not necessarily in the admissible dual).  Given $\phi$ as above, we may choose a compact open subgroup $U$ of $G_n$ under which $\phi$ is bi-invariant.  The integrals in question do not change if we then replace $f$ by its two-sided average with respect to $U$, which has the effect of replacing $v$ by its average $v^U \in \pi^U$ and $\ell$ with its projection $\ell^U$ to the dual of $\pi^U$, extended by zero on the kernel of the averaging operator $\pi \rightarrow \pi^U$.  In particular, by specializing to the case that $\ell$ is a Whittaker functional on $\pi$, we see that such identities remain valid when $f$ is a Whittaker function for $\pi$.

We denote by $\mathcal{S}^{e}(F^\times)$ the space of all Schwartz--Bruhat functions $\beta \in \mathcal{S}(F^\times)$ such that $\beta(x y) = \beta(x)$ whenever $\lvert y \rvert = 1$, or equivalently, for which $\beta(x)$ depends only upon $\lvert x\rvert$.  We note that each $\beta \in \mathcal{S}^{e}(F^\times)$ satisfies the Mellin inversion formula
\begin{equation}\label{eq:cq73uvecsz}
  \beta(y) = \int_{(\sigma)} \tilde{\beta}(s) \lvert y \rvert^s
  \, d s,
  \qquad \tilde{\beta}(s) := \int_{F^\times} \beta(y) \lvert y \rvert^{-s} \,d^\times y.
\end{equation}

For $\beta \in \mathcal{S}^e(F^\times)$, we define the transform $\beta^\sharp := \beta^{\sharp, \pi}$ of $\beta$ by
\begin{equation*}
    \beta^\sharp(y) :=
    \int_{(\sigma)}
    \frac{
      \tilde{\beta}(s)
      \lvert y \rvert^{- s}
      \, d s}{\gamma(\tfrac{1}{2} + s, \pi, \psi)},
  \end{equation*}
initially for $\sigma$ large enough.
\begin{lemma}\label{example:cq4jqru7rh}
  Define $\beta$ via Mellin inversion \eqref{eq:cq73uvecsz} by
  \begin{equation*}
    \tilde{\beta}(s) := \frac{\varepsilon(\tfrac{1}{2} + s, \pi, \psi)}{L(\tfrac{1}{2} + s, \pi)}.
  \end{equation*}
  Then:
  \begin{enumerate}
  \item $\beta$ is supported on $\{y : \lvert Q \rvert \leq \lvert y \rvert \leq \lvert Q \rvert q^n \}$ and takes the value $\varepsilon(\tfrac{1}{2}, \pi, \psi)$ on $\{y : \lvert y \rvert = \lvert Q \rvert\}$.
  \item $\beta^\sharp$ is supported on $\{y : 1 \leq \lvert y \rvert \leq q^n\}$ and takes the value $1$ on $\{y : \lvert y \rvert = 1\} = \mathfrak{o}^\times$.
  \end{enumerate}
\end{lemma}
\begin{proof}
  We appeal to the characterization \eqref{eq:cuipy1wony} of $\lvert Q \rvert$.  We note first that $\beta^\sharp$ has Mellin transform
  \begin{equation*}
    \widetilde{\beta^\sharp} (s) = \frac{1}{L(\tfrac{1}{2} + s, \tilde{\pi})}.
  \end{equation*}
  Since the inverse $L$-values appearing above are monic polynomials in $q^{- s}$ of degree at most $n$, we see by Mellin inversion that $\beta$ and $\beta^\sharp$ have the claimed properties.
\end{proof}

\begin{lemma}\label{lemma:cq4itbu91o}
  Assume that $\pi$ is unitary and generic.  We then have the identity of absolutely convergent integrals
  \begin{equation}\label{eq:cq4jqjygcl}
    \int_{G_n}
    \phi(g)
    f(g)
    \beta(\det g)
    \lvert \det g \rvert^{\frac{n}{2}}
    \, d g
    =
    \int_{G_n}
    \phi^\wedge(g)
    f^\vee(g)
    \beta^\sharp(\det g) \lvert \det g \rvert^{\frac{n}{2}} \, d g.
  \end{equation}
\end{lemma}
\begin{proof}
  Starting with the left hand side, we insert the Mellin expansion of $\beta$, with $\sigma = 0$.  The resulting double integral over $g$ and $s$ converges absolutely, so we may swap the order.  We recognize the result as the integral $\int_{(0)} \tilde{\beta}(s) Z(\phi, f, \tfrac{1}{2} + s) \, d s$ involving the Godement--Jacquet zeta integral \eqref{eq:cq73uvikdb}.  We now apply the local functional equation and expand the result as
  \begin{equation*}
    \int_{(0)}
    \frac{\tilde{\beta}(s)}{\gamma(\tfrac{1}{2} + s, \pi, \psi)}
    \left(
      \int_{G_n}
      \phi^\wedge(g) f^\vee(g) \lvert \det g \rvert^{\frac{n}{2} - s}
      \, d g
    \right)
    \, d s.
  \end{equation*}
  This double integral again converges absolutely, so we may rearrange it to obtain the stated identity.
\end{proof}
For the same reasons as indicated following the statement of Lemma \ref{lemma:cq4itap0v1}, such identities persist for more general coefficients than matrix coefficients, and in particular, when $f$ is a Whittaker function.


Recall that we embed $G_n \hookrightarrow G_{n + 1}$ as the upper-left block.  We set
\begin{equation}\label{eq:cq4jqk9bz0}
  W_0(g):=\int_{N_n} 1_{K_n}(x g)\,\psi(x)\,d x,
\end{equation}
which defines a Whittaker function on $G_n$ and extends, by the theory of the Kirillov model \cite{MR748505}, to an element of $\mathcal{W}(\Pi,\psi^{-1})$ on $G_{n+1}$.

For $x \in F$ and $y \in F^\times$, we set
\begin{equation*}
  d_y:=\operatorname{diag}(1,\dotsc,1,y)\in G_n \hookrightarrow G_{n + 1},
  \qquad
  u_x:=I_{n+1}+ x E_{n,n+1} \in N_{n + 1}.
\end{equation*}
We then define
\begin{equation*}
  t_Q:= d_Q^{-1} u_Q = u_1 d_Q^{-1}.
\end{equation*}
\begin{lemma}\label{proposition:cq4jqgh7yi}
  There exist $\beta \in \mathcal{S}^e(F^\times)$ and $\phi \in \mathcal{S}(M_n(F))$ so that for all $g \in G_n$, we have
  \begin{equation}\label{eq:cq4jqjeuyf}
    \int_{N_n} \beta(\det x g) \phi(x g) \psi(x) \, d x = \varepsilon(\tfrac{1}{2}, \pi, \psi) W_0(g t_{Q})
  \end{equation}
  and
  \begin{equation}\label{eq:cq4jqjy05j}
    \beta^\sharp(\det g) \phi^\wedge(g) = \lvert Q \rvert^n 1_{K_1(\mathfrak{q})}(g).
  \end{equation}
\end{lemma}

\begin{proof}
  We set
  \begin{equation*}
    \phi_0 := 1_{M_n(\mathfrak{o})},
  \end{equation*}
  \begin{equation}\label{eq:cq4jqv74tm}
    \phi(x) := \psi(- x_{nn}) \phi_0(x d_{Q}^{-1}).
  \end{equation}
  and take $\beta$ as in Lemma \ref{example:cq4jqru7rh}, so that in particular,
  \begin{equation}\label{eq:cq73v9ep9j}
    \beta|_{Q \mathfrak{o}} = \varepsilon(\tfrac{1}{2}, \pi, \psi) 1_{Q \mathfrak{o}^\times}
  \end{equation}
  and
  \begin{equation}\label{eq:cq73v91fxa}
    \beta^\sharp |_{\mathfrak{o}} = 1_{\mathfrak{o}^\times}.
  \end{equation}
  We must verify the relations \eqref{eq:cq4jqjeuyf} and \eqref{eq:cq4jqjy05j}.

  We start with \eqref{eq:cq4jqjeuyf}.  Recall from \eqref{eq:cq4jqk9bz0} that $W_0$ is the $\psi^{-1}$-Whittaker function $W_0(g) = \int_{N_n} 1_{K_n}(x g) \psi(x) \, d x$.  In particular,
  \begin{equation}\label{eq:cq4jqk3rcx}
    W_0(g t_{Q}) = W_0(g u_1 d_{Q}^{-1}) = \psi(-g_{n n}) W_0(g d_{Q}^{-1}).
  \end{equation}
  Using this identity, we may rewrite the desired relation \eqref{eq:cq4jqjeuyf} as
  \begin{equation}\label{eq:cq6s1gjyh4}
    \int_{N_n} \beta(\det (x g)) \phi(x g) \psi(x) \, d x = \varepsilon(\tfrac{1}{2}, \pi, \psi) \psi(- g_{n n}) W_0(g d_{Q}^{-1}).
  \end{equation}
  We verify this as follows.  First, we see from the definition \eqref{eq:cq4jqv74tm} and the identity $(xg)_{nn} = g_{nn}$ that for $x \in N_n$ and $g \in G_n$, we have
  \begin{equation}\label{eq:relation-phi-phi0}
    \phi(x g) = \psi(- g_{n n}) \phi_0(x g d_Q^{-1}).
  \end{equation}
  Next, we have
  \begin{align*}
    \beta(\det g) \phi_0(g d_Q^{-1}) &= \varepsilon(\tfrac{1}{2}, \pi, \psi)
                                1_{Q \mathfrak{o}^\times}(\det g)
                                \phi_0(g d_Q^{-1}) \\
                              &=
                                \varepsilon(\tfrac{1}{2}, \pi, \psi)
                                1_{K_n}(g d_Q^{-1}).
  \end{align*}
  (In the first step, we use that $\phi_0(g d_Q^{-1})$ is nonzero only if $\det(g) \in Q \mathfrak{o}$ and apply \eqref{eq:cq73v9ep9j}.  In the second step, we use that $1_{K_n}(g) = 1_{\mathfrak{o}^\times}(\det g) \phi_0(g)$ and $\det(d_Q) = Q$, which gives $1_{Q \mathfrak{o}^\times}(\det g) \phi_0(g d_Q^{-1}) = 1_{K_n}(g d_Q^{-1})$.)  Combining the above identities, we obtain
  \begin{equation*}
    \beta(\det(x g)) \phi(x g) = \varepsilon(\tfrac{1}{2}, \pi, \psi) \psi(- g_{n n}) 1_{K_n}(x g d_Q^{-1}).
  \end{equation*}
  Integrating both sides against $\psi(x) \, d x$ gives \eqref{eq:cq6s1gjyh4}, as required.

  We verify \eqref{eq:cq4jqjy05j} as follows (here $E_{ij}$ denotes the elementary matrix):
  \begin{align*}
    \beta^\sharp(\det g) \phi^\wedge(g)
    &=
      1_{\mathfrak{o}^\times}(\det g) \phi^\wedge(g)\\
    &=
      1_{\mathfrak{o}^\times}(\det g)
      \lvert Q \rvert^n \phi_0^\wedge(d_{Q}(g - E_{nn})) \\
    &=
      \lvert Q \rvert^n
      1_{\mathfrak{o}^\times}(\det g)
      1_{M_n(\mathfrak{o})}(d_{Q}(g - E_{nn})) \\
    &=
      \lvert Q \rvert^n
      1_{K_1(\mathfrak{q})}(g).
  \end{align*}
  Here, for the first step, we observed that $\phi^\wedge(x)$ is nonzero only if $x \in E_{nn} + d_Q^{-1} M_n(\mathfrak{o}) \subseteq M_n(\mathfrak{o})$, so that, in particular, $\det x \in \mathfrak{o}$; we then applied \eqref{eq:cq73v91fxa}.  For the second step, we applied the general Fourier analytic calculation
  \begin{equation}\label{eq:cq4jqkb0q8}
    \phi^\wedge(x) = \lvert Q \rvert^n \phi_0^\wedge(d_{Q}(x - E_{nn})).
  \end{equation}
  For the third, we applied the Fourier self-duality $\phi_0^\wedge = \phi_0 = 1_{M_n(\mathfrak{o})}$.  For the final step, we use that $K_1(\mathfrak{q})$ consists of all $x \in M_n(F)$ for which $d_Q (x - E_{nn}) \in M_n(\mathfrak{o})$ and $\det x \in \mathfrak{o}^\times$.
\end{proof}

For $W \in \mathcal{W}(\Pi, \psi^{-1})$, $V \in \mathcal{W}(\pi, \psi)$, and $s \in \mathbb{C}$, we define the Rankin--Selberg integral
\begin{equation}\label{eq:cq4fl8vokh}
  \ell_{\mathrm{RS}}(s, W, V) := \int_{N_n \backslash G_n} W (\operatorname{diag}(g, 1)) \, V(g) \, \lvert \det g \rvert^{s - \frac{1}{2}} \, d g.
\end{equation}
The following result verifies Theorem \ref{thm:main} in a more precise form.
\begin{proposition}\label{proposition:cq4jqgikw1}
  Let $W_0 \in \mathcal{W}(\Pi, \psi^{-1})$ be such that for all $g \in G_n$, we have
  \begin{equation*}
    W_0(g) = \int_{N_n} 1_{K_n}(x g) \psi(x) \, d x.
  \end{equation*}
  Let $V \in \mathcal{W}(\pi, \psi)$ denote the normalized newvector (i.e., the unique $K_1(\mathfrak{q})$-invariant vector for which $V(1) = 1$, see \cite{MR620708, MR3138844}).  Then for all $s \in \mathbb{C}$, we have
  \begin{equation}\label{eq:cq4jqi2rfd}
    \ell_{\mathrm{RS}}(s, u_Q W_0, d_Q V) = c \lvert Q \rvert^{- \frac{n}{2}},
  \end{equation}
  where
  \begin{equation}\label{eq:cq4jqlhqy9}
    c := \varepsilon(\tfrac{1}{2}, \pi, \psi)^{-1} \lvert Q \rvert^n \operatorname{vol}(K_1(\mathfrak{q})) \asymp 1.
  \end{equation}
\end{proposition}
\begin{proof}
  We note first that, by a change of variables, we have the homogeneity property
  \begin{equation}\label{eq:cq4jqpwhnv}
    \ell_{\mathrm{RS}}(s, u_{Q} W_0, d_{Q}V) = \lvert Q \rvert^{- \left( s - \frac{1}{2} \right)} \ell_{\mathrm{RS}}(s, t_{Q} W_0, V).
  \end{equation}
  In view of this, the desired identity \eqref{eq:cq4jqi2rfd} is equivalent to
  \begin{equation}\label{eq:cq4jqjdmzj}
    \ell_{\mathrm{RS}}(s, t_{Q} W_0, V) = c \lvert Q \rvert^{s - \frac{n+1}{2}}.
  \end{equation}
  Next, since $W_0$ is supported on $\det^{-1}(\mathfrak{o}^\times)$, we see that the translate $t_{Q} W_0$ is supported on $\det^{-1}(Q \mathfrak{o}^\times)$, so the left hand side of \eqref{eq:cq4jqjdmzj} is a constant multiple of $\lvert Q \rvert^s$.  For this reason, it suffices to verify \eqref{eq:cq4jqjdmzj} for (say) $s = \tfrac{n+1}{2}$, where our task is to check that $\ell_{\mathrm{RS}}(\tfrac{n+1}{2}, t_{Q} W_0, V) = c$.
  Inserting definitions and unfolding, we obtain, with $f(g) := V(g)$,
  \begin{align*}
    \varepsilon(\tfrac{1}{2}, \pi, \psi) \ell_{\mathrm{RS}}(\tfrac{n+1}{2}, t_{Q} W_0, V)
    &\overset{\eqref{eq:cq4fl8vokh}}{=}
      \varepsilon(\tfrac{1}{2}, \pi, \psi)
      \int_{N_n \backslash G_n} W_0(g t_Q) V (g) \lvert \det (g) \rvert^{\frac{n}{2}} \, d g
    \\
    &\overset{\eqref{eq:cq4jqjeuyf}}{=}
      \int_{G_n}
      \phi(g)
      f(g)
      \beta(\det g)
      \lvert \det g \rvert^{n/2}
      \, d g \\
    &\overset{\eqref{eq:cq4jqjygcl}}{=}
      \int_{G_n}
      \phi^\wedge(g)
      f^\vee(g)
      \beta^\sharp(\det g) \lvert \det g \rvert^{n/2} \, d g \\
    &\overset{\eqref{eq:cq4jqjy05j}}{=}
      \lvert Q \rvert^n \int_{K_1(\mathfrak{q})} V(g^{-1}) \lvert \det g \rvert^{n/2} \, d g \\
    &= \lvert Q \rvert^n \operatorname{vol}(K_1(\mathfrak{q})),
  \end{align*}
  where in the final step, we use the $K_1(\mathfrak{q})$-invariance of $V$, the normalization $V(1) = 1$, and the fact that $\lvert \det g \rvert = 1$ on $K_1(\mathfrak{q})$.  Thus \eqref{eq:cq4jqjdmzj} holds.
\end{proof}

% \bibliography{refs}{} \bibliographystyle{plain}

\def\cprime{$'$} \def\cprime{$'$} \def\cprime{$'$} \def\cprime{$'$}
\begin{thebibliography}{1}

\bibitem{MR748505}
  Joseph~N. Bernstein.
  \newblock {$P$}-invariant distributions on {${\rm GL}(N)$} and the
  classification of unitary representations of {${\rm GL}(N)$}
  (non-{A}rchimedean case).
  \newblock In {\em Lie group representations, {II} ({C}ollege {P}ark, {M}d.,
    1982/1983)}, volume 1041 of {\em Lecture Notes in Math.}, pages 50--102.
  Springer, Berlin, 1984.

\bibitem{MR4053059}
  Andrew~R. Booker, M.~Krishnamurthy, and Min Lee.
  \newblock Test vectors for {R}ankin-{S}elberg {$L$}-functions.
  \newblock {\em J. Number Theory}, 209:37--48, 2020.

\bibitem{MR0342495}
  Roger Godement and Herv{\'e} Jacquet.
  \newblock {\em Zeta functions of simple algebras}.
  \newblock Lecture Notes in Mathematics, Vol. 260. Springer-Verlag, Berlin,
  1972.

\bibitem{MR620708}
  H.~Jacquet, I.~I. Piatetski-Shapiro, and J.~Shalika.
  \newblock Conducteur des repr\'esentations du groupe lin\'eaire.
  \newblock {\em Math. Ann.}, 256(2):199--214, 1981.

\bibitem{MR3138844}
  Nadir Matringe.
  \newblock Essential {W}hittaker functions for {$GL(n)$}.
  \newblock {\em Doc. Math.}, 18:1191--1214, 2013.

\bibitem{michel-2009}
  Philippe Michel and Akshay Venkatesh.
  \newblock The subconvexity problem for {${\rm GL}_2$}.
  \newblock {\em Publ. Math. Inst. Hautes \'Etudes Sci.}, (111):171--271, 2010.

\bibitem{MR780071}
  Peter Sarnak.
  \newblock Fourth moments of {G}r\"ossencharakteren zeta functions.
  \newblock {\em Comm. Pure Appl. Math.}, 38(2):167--178, 1985.

\end{thebibliography}

\end{document}
